\chapter{Introduction}

Les réseaux de neurones convolutionnels, souvent appelés CNN pour \textit{Convolutional Neural Networks}, 
inspirés par le cortex visuel des vertébrés, occupent une place centrale en vision par ordinateur. 
Ils sont utilisés pour reconnaître des objets, classer des images, détecter des formes ou encore extraire 
automatiquement des informations visuelles. 

Leur efficacité vient en grande partie de leur capacité à apprendre des motifs locaux dans les 
images, puis à les combiner progressivement pour construire des représentations plus abstraites. 

\section*{Fragilités et attaques adversariales}

Cependant, malgré leurs très bonnes performances, ces modèles peuvent présenter des fragilités structurelles. 
Une image peut être faiblement modifiée, de telle manière que cela devienne imperceptible à l'oeil humain, 
tout en entraînant un changement important dans la prédiction du réseau. Ce phénomène est connu sous le nom de 
\textbf{perturbation adversariale}~\cite{szegedy2014intriguing}.

Un aspect particulièrement intéressant de ces perturbations est leur caractère parfois universel. Dans certains cas, 
il est possible de trouver une perturbation qui est directement basée sur l'architecture du réseau 
(sa manière de fonctionner), exploitant une faiblesse structurelle du fonctionnement interne de ce dernier. 
Ainsi ce genre de perturbation, une fois trouvée, peut être appliquée à n'importe quelle image, 
\textbf{indépendamment} de son contenu, et tromper le modèle, pour une proportion importante d'images,
dans ses prédictions. On parle alors de perturbation adversariale \textbf{universelle}~\cite{moosavi2017universal}. 

La question centrale est alors de comprendre d'où vient cette sensibilité. Une première explication possible 
consiste à l'attribuer aux données d'entraînement ou à l'optimisation du modèle. Dans ce travail, nous nous 
intéressons à une autre piste : \textbf{la structure même des réseaux convolutionnels}. En effet, les CNN 
reposent sur des couches de convolution, et ces couches possèdent des propriétés mathématiques particulières. 

\section*{L'approche par la base de Fourier}

Pour étudier cette question, nous choisissons d'utiliser la base de Fourier. Une image peut être 
décrite pixel par pixel, mais elle peut aussi être décomposée en motifs fréquentiels : des variations lentes, 
des oscillations rapides, des motifs horizontaux, verticaux ou diagonaux. Ces motifs fréquentiels donnent une 
autre manière de décrire les structures présentes dans une image.

Ainsi, analyser une image non plus seulement dans le domaine spatial, c'est-à-dire à partir de ses pixels, mais 
aussi dans le domaine fréquentiel, permet de mieux comprendre quelles composantes de l'image influencent le 
comportement du réseau et quelles directions de perturbation peuvent devenir particulièrement sensibles.

Ce lien constitue le point de départ du travail de Tsuzuku et Sato~\cite{tsuzuku2019structural}, qui montrent 
que certains réseaux convolutionnels sont particulièrement sensibles à des directions issues de la base de Fourier. 
Leur idée est que certaines perturbations adversariales universelles peuvent être mieux comprises en étudiant la 
manière dont les couches convolutionnelles traitent les fréquences.

\begin{center}
    \vspace{1em}
    \rule{0.4\textwidth}{0.4pt} \\
    \vspace{1em}
    \begin{minipage}{0.85\textwidth}
        \centering
        \large \itshape
        Comment l'architecture convolutionnelle des CNN induit-elle des directions de vulnérabilité privilégiées dans 
        le domaine de Fourier, et dans quelle mesure cette sensibilité structurelle permet-elle de mieux comprendre 
        l'existence de perturbations adversariales universelles ?
    \end{minipage} \\
    \vspace{1em}
    \rule{0.4\textwidth}{0.4pt}
    \vspace{1em}
\end{center}

L'objectif de ce rapport est de reprendre cette idée et de la reconstruire petit à petit de manière plus détaillée en 
introduisant les éléments nécessaires pour comprendre les résultats, de telle manière que n'importe quel lecteur avec une connaissance 
de base en algèbre linéaire et en machine learning puisse comprendre les implications de ce travail.

\section*{Organisation du rapport}

Pour répondre à cette problématique, le rapport est organisé de la manière suivante :

\begin{itemize}[label=\textbullet, leftmargin=1.5cm]
    \item \textbf{Chapitre 2 :} Présentation du contexte général : réseaux convolutionnels, robustesse et perturbations adversariales.
    \item \textbf{Chapitre 3 :} Présentation des fondements théoriques reliant convolution et base de Fourier. 
    On y montre notamment l'émergence naturelle des fonctions de Fourier dans l'analyse linéaire.
    \item \textbf{Chapitre 4 :} Présentation de la reproduction expérimentale, du protocole utilisé, des \textit{heatmaps} 
    de sensibilité fréquentielle, puis discussion des résultats et des limites.
    \item \textbf{Conclusion :} Synthèse des enseignements et perspectives de recherche.
\end{itemize}
